% 可数集（综述）
% license CCBYSA3
% type Wiki

本文根据 CC-BY-SA 协议转载翻译自维基百科\href{https://en.wikipedia.org/wiki/Countable_set}{相关文章}

在数学中，如果一个集合是有限的，或者它可以与自然数集合建立一一对应关系，那么它就是可数的。[a] 等价地说，如果存在一个从该集合到自然数的单射函数，那么该集合就是可数的；这意味着集合中的每个元素都可以与一个唯一的自然数对应，或者说集合的元素可以一个接一个地数出来，尽管由于元素个数无限，计数过程可能永远不会结束。

从更技术性的角度讲，假设接受可数选择公理，若一个集合的基数（即集合中的元素个数）不大于自然数集合的基数，则它是可数的。一个可数但不是有限的集合称为可数无穷。

这一概念归功于格奥尔格·康托尔，他证明了不可数集的存在，也就是说存在一些集合不是可数的；例如实数集合就是不可数的。
\subsection{关于术语的说明}
尽管这里所定义的“可数”与“可数无穷”用法相当普遍，但并非所有文献都采用这种术语。[1] 一种替代用法是：用countable 来表示这里所说的“可数无穷”，而用at most countable来表示这里所说的“可数”。[2][3]

术语enumerable[4]和denumerable[5][6] 也可能被使用，例如分别指代“可数”和“可数无穷”。[7] 但其定义并不统一，并且需要注意与“递归可枚举”的区别。[8]
\subsection{定义}
一个集合$S$是可数的，当且仅当满足以下任一条件：
\begin{itemize}
\item 它的基数 $|S|$ 小于或等于 $\aleph_{0}$（阿列夫零，$\mathbb{N}$ ——自然数集的基数）。[9]
\item 存在一个从 $S$ 到 $\mathbb{N}$ 的单射函数。[10][11]
\item $S$ 是空集，或者存在一个从 $\mathbb{N}$ 到 $S$ 的满射函数。[11]
\item 存在 $S$ 与 $\mathbb{N}$ 的某个子集之间的双射映射。[12]
\item $S$ 要么是有限的（$|S| < \aleph_{0}$），要么是可数无穷的。[5]
\end{itemize}
以上定义都是等价的。

一个集合 $S$ 是**可数无穷的**，当且仅当：
\begin{itemize}
\item 它的基数 $|S|$ 正好等于 $\aleph_{0}$。[9]
\item 存在一个从 $S$ 到 $\mathbb{N}$ 的既单射又满射（因此是双射）的映射。
\item $S$ 与 $\mathbb{N}$ 之间存在一一对应关系。[13]
\item $S$ 的元素可以排列成一个无限序列$a_{0}, a_{1}, a_{2}, \ldots$其中 $a_i \neq a_j$（当 $i \neq j$ 时），并且 $S$ 的每个元素都在该序列中出现。[14][15]
\end{itemize}
如果一个集合不是可数的，即它的基数大于 $\aleph_{0}$，那么它就是不可数的。[9]
\subsection{历史}
1874 年，康托尔在他的第一篇集合论文章中证明了实数集是不可数的，从而表明并非所有无限集合都是可数的。[16]1878 年，他利用一一对应来定义并比较基数。[17]1883 年，他用无限序数扩展了自然数，并利用序数集合构造出具有不同无限基数的无限多个集合。[18]
\subsection{引言}
集合是元素的一个整体，可以通过多种方式来描述。一种方式就是直接列出所有元素；例如，包含整数 3、4 和 5 的集合可以记作$\{3, 4, 5\}$，这称为列举法。[19]然而，这种方法只对小集合有效；对于较大的集合，这会非常耗时且容易出错。为了避免逐一列出所有元素，有时会使用省略号“...”来表示从起始元素到结束元素之间的很多元素，前提是作者认为读者能很容易推测出“...”所代表的内容。例如：$\{1,2,3,\dots,100\}
$通常表示从 1 到 100 的所有整数。即便在这种情况下，仍然有可能把所有元素完全列出，因为这个集合中的元素个数是有限的。如果我们依次给集合中的元素编号为 1，2，…，直到$n$，这就给出了“大小为$n$的集合”的通常定义。
\begin{figure}[ht]
\centering
\includegraphics[width=6cm]{./figures/009b7b48421ad40c.png}
\caption{} \label{fig_KSJ_1}
\end{figure}
有些集合是无限的；这些集合的元素数目大于任意一个可以指定的整数$n$。 （无论指定的整数$n$多大，比如 $n=10^{1000}$，无限集中的元素都比$n$多。）例如，自然数集合可以记作：$\{0,1,2,3,4,5,\dots\}$，它有无穷多个元素，我们无法用任何一个自然数来表示它的大小。[a]看起来很自然的做法是将集合分为不同的类：把所有只含有一个元素的集合放在一起；所有含有两个元素的集合放在一起；……；最后，把所有无限集合放在一起，并认为它们大小相同。这种观点对于可数无限集是成立的，并且在格奥尔格·康托尔之前是普遍的假设。例如，奇数集合是无限的，偶数集合是无限的，整数集合也是无限的。我们可以认为这些集合大小相同，因为我们可以这样安排：对每一个整数，都对应一个唯一的偶数：
$$
\ldots,\,-2 \mapsto -4,\,-1 \mapsto -2,\,0 \mapsto 0,\,1 \mapsto 2,\,2 \mapsto 4,\ldots~
$$
更一般地：$n \mapsto 2n$（见图）。这里我们所做的是在整数集合与偶数集合之间建立了一一对应（双射）。双射是两个集合之间的一个映射，使得每个集合中的每个元素都唯一对应于另一个集合中的某个元素。这种数学上的“大小”概念称为基数：两个集合大小相同，当且仅当它们之间存在双射。所有与整数集合存在一一对应的集合称为可数无穷集，并且它们的基数记作 $\aleph_{0}$。

康托尔证明了并非所有无限集都是可数无穷的。例如，实数集合不能与自然数集合（一切非负整数）建立一一对应关系。实数集的基数大于自然数集，因此称为不可数的。
\subsection{形式化概述}
根据定义，一个集合 $S$ 是**可数的**，如果存在一个 $S$ 与自然数集合 $\mathbb{N}=\{0,1,2,\dots\}$ 的某个子集之间的双射。

例如，定义对应关系：
$$
a \leftrightarrow 1, \quad b \leftrightarrow 2, \quad c \leftrightarrow 3~
$$
由于集合$S=\{a,b,c\}$中的每个元素都恰好与集合 $\{1,2,3\}$ 中的一个元素配对，反之亦然，这就构成了一个双射，并说明 $S$ 是可数的。类似地，我们可以证明所有有限集合都是可数的。

对于无限集合的情况：一个集合$S$是可数无穷的，当且仅当存在一个$S$与整个自然数集合$\mathbb{N}$之间的双射。例如，考虑集合$A = \{1,2,3,\dots\}$（正整数集合）和$B = \{0,2,4,6,\dots\}$（偶整数集合）。我们可以通过展示它们与自然数集的双射来说明这些集合是可数无穷的。这可以通过如下映射实现：$n\leftrightarrow n+1, \quad n \leftrightarrow 2n$于是：
$$
\begin{matrix}  
0 \leftrightarrow 1, & 1 \leftrightarrow 2, & 2 \leftrightarrow 3, & 3 \leftrightarrow 4, & 4 \leftrightarrow 5, & \dots \\[6pt]  
0 \leftrightarrow 0, & 1 \leftrightarrow 2, & 2 \leftrightarrow 4, & 3 \leftrightarrow 6, & 4 \leftrightarrow 8, & \dots  
\end{matrix} ~
$$
这就分别建立了自然数集与正整数集、自然数集与偶整数集之间的一一对应关系，从而说明它们是可数无穷的。

每一个可数无穷集都是可数的，而每一个无限可数集也是可数无穷的。此外，自然数集合的任意子集都是可数的，更一般地：

\textbf{定理} —— 可数集的子集仍是可数的。[20]

所有自然数有序对的集合（即自然数集合的笛卡尔积$\mathbb{N} \times \mathbb{N}$）是可数无穷的，这一点可以通过按照图中所示路径进行枚举来说明：

得到的映射过程如下：
\begin{figure}[ht]
\centering
\includegraphics[width=6cm]{./figures/16cb3ab4d89837dc.png}
\caption{} \label{fig_KSJ_2}
\end{figure}
$$
0 \leftrightarrow (0,0),\; 
1 \leftrightarrow (1,0),\; 
2 \leftrightarrow (0,1),\; 
3 \leftrightarrow (2,0),\; 
4 \leftrightarrow (1,1),\; 
5 \leftrightarrow (0,2),\; 
6 \leftrightarrow (3,0),\; 
\ldots~
$$
这种映射覆盖了所有这样的有序对。

这种三角形映射的形式可以递归推广到自然数的 $n$-元组，即$(a_{1}, a_{2}, a_{3}, \dots , a_{n})$，其中 $a_i$ 和 $n$ 都是自然数。方法是不断地将 $n$-元组的前两个元素映射为一个自然数。例如，$(0,2,3)$可以写作$((0,2),3)$。其中，$(0,2)$ 映射为 5，因此 $((0,2),3)$ 映射为 $(5,3)$，然后 $(5,3)$ 映射为 39。由于不同的二元组（例如 $(a,b)$）会映射到不同的自然数，因此只要两个 $n$-元组在某一元素上不同，就能保证它们被映射到不同的自然数。由此证明了从所有 $n$-元组集合到自然数集合 $\mathbb{N}$ 的一个单射存在。对于由有限多个不同集合的笛卡尔积构成的 $n$-元组，每个元组中的每个元素都可以对应到一个自然数，因此每个元组都可以写成自然数形式，然后应用同样的逻辑即可证明该结论。

\textbf{定理}—— 有限多个可数集合的笛卡尔积仍然是可数的。[21][b]

